% !TeX root = ../../main.tex
\subsection{Binary padding that hides two extension-field evaluations}\label{note:zk-two-point-span}

\paragraph{Purpose.} Counter-label swaps supply fair bits, not uniform extension-field values. To preserve entropy after an earlier bus evaluation, one needs a \emph{joint} two-point span theorem. The single-point lemma cannot be applied twice to the same bits. The construction below gives a uniform statistical bound using $576$ bits; no computational pseudorandomness is assumed.

Let $Q=|\E|=2^{192}$ and let $S\subset\cube{14}$ have arbitrary first six bits and Hamming weight at most one on the last eight. Thus $|S|=64\cdot9=576$.

\begin{theorem}\label{thm:two-point-span}
For independent uniform $r,u\in\E^{14}$, the vectors
\[
 (\eq(r,s),\eq(u,s)),\qquad s\in S,
\]
span $\E^2/\mathbb F_2$ except with probability below $2^{-148}$. The same bound holds after multiplying both coordinates at position $s$ by $\prod_i\lambda_i^{s_i}$, for arbitrary fixed nonzero $\lambda_i\in\K$. Independent fair bits at these positions therefore give a uniform pair of evaluations on good coins, not merely two uniform marginals.
\end{theorem}

\subsubsection{A product-space span argument}

First work with normalized monomial weights and independent scalars in $\E^*$. The first six coordinates generate a binary subspace $V\subset\E^2$. Its two projected dimensions are $d_1,d_2$ and its full dimension is $d_{12}\ge\max(d_1,d_2)$. The last eight coordinates give eight independent random diagonal multiples of $V$, in addition to $V$ itself.

Use the nondegenerate trace pairing on $\E^2$. For a nonempty coordinate support $I\subseteq\{1,2\}$, at most $2^{192|I|-d_I}$ dual vectors supported there annihilate $V$. A fixed such nonzero vector annihilates one random diagonal multiple with probability at most
\[
 \frac{2^{192|I|-d_I}}{(Q-1)^{|I|}}.
\]
After eight independent multiples, a union bound gives the support-$I$ failure bound
\[
 \left(\frac Q{Q-1}\right)^{8|I|}2^{192|I|-9d_I}
 \le 2^{192|I|+1-9d_I}.
\]
Cap each probability at one. The two singleton supports are controlled by their own projected dimensions; the two-coordinate support can conservatively use $d_{12}\ge d_1$.

\subsubsection{Averaging the sixth doubling}

For each projection, the first five doublings reach dimension $32$ except with probability
\[
 B=\sum_{d\in\{1,2,4,8,16\}}\frac{(2^d-1)^2}{Q-1}.
\]
Indeed, a bad multiplier is a ratio of two nonzero elements of the current subspace. Condition on both projections reaching dimension $32$, costing at most $2B$. For an arbitrary fixed $32$-dimensional subspace $W\subset\E$, let $D=\dim(W+tW)$ after the sixth scalar $t\gets\E^*$. Then $32\le D\le64$, and for $32\le d<64$,
\begin{equation}\label{eq:two-dimension-tail}
 \Pr[D=d]\le\Pr[D\le d]
 \le p_d:=\frac{(2^{32}-1)^2}{(Q-1)(2^{64-d}-1)}.
\end{equation}
To see this, an intersection of dimension at least $64-d$ contains at least $2^{64-d}-1$ nonzero pairs satisfying $v=tw$. Across all $t$, there are only $(2^{32}-1)^2$ such pairs. This tail bound is sharper than discarding every failed sixth doubling, which would lose the required statistical margin.

Define
\[
 A_a=\min\{1,2^{a+1-9\cdot64}\}
 +\sum_{d=32}^{63}p_d\min\{1,2^{a+1-9d}\}.
\]
The normalized two-point failure probability is at most
\[
 B_2=2B+2A_{192}+A_{384}.
\]
No independence between projected dimensions is needed: the argument bounds each singleton term separately and uses just one projection for the last term.

\subsubsection{Returning to equality weights}

For each of the $28$ point coordinates, exclude $r_i=1$ and normalize with $t_i=r_i/(1+r_i)$. Excluding zero normalized values costs at most $28/(Q-1)$; their remaining distribution differs from uniform $\E^*$ by at most another $28/(Q-1)$ in total. Nonzero common factors on the two output coordinates do not affect rank. Thus ordinary equality weights have failure at most
\[
 B_2+\frac{28}{Q}+\frac{56}{Q-1}.
\]
For geometric multipliers, use the coordinate substitution from \noteref{zk-jump-metadata}, simultaneously at the two points. Poles and the distance of the transformed points from uniform contribute at most another $56/Q$. Exact rational arithmetic verifies
\[
 B_2+\frac{84}{Q}+\frac{56}{Q-1}<2^{-148}.
\]

\paragraph{Evidence and boundary.} \auditref{zk_two_point_audit.py} checks this rational inequality and binary rank $384$ for ordinary and geometrically weighted evaluations over the actual fields. The theorem requires two independent honest field points. It does not automatically cover a third query, an adaptive earlier transcript or arbitrary PCS selectors. Its intended first use is a counter bank whose lower-half swaps affect two evaluations while leaving a separately masked upper-half bus target unchanged.

\subsubsection{A compact construction within the PCS size cap}\label{sec:zk-compact-two-point}

The preceding construction uses fourteen coordinates. For Flock, a denser support on eleven coordinates leaves room for the rest of the committed witness. Define
\[
 \zkCompactSupport=\{s\in\cube{11}:s_0,\ldots,s_6\text{ arbitrary},\quad s_7+\cdots+s_{10}\le1\},
 \qquad |\zkCompactSupport|=640,
\]
where the inequality is over integers. The improvement uses the dimension of the \emph{joint} two-coordinate space, rather than bounding it by one projected dimension.

\begin{theorem}[Compact two-point span]\label{thm:zk-compact-two-point}
For two independent uniform points in $\E^{11}$, the paired equality weights on $\zkCompactSupport$ span $\E^2$ over $\Ftwo$ except with probability $\zkCompactError<2^{-158}$. The same upper bound holds with fixed nonzero geometric multipliers as in Theorem~\ref{thm:two-point-span}.
\end{theorem}
\begin{proof}
Work first with independent nonzero monomial multipliers, and let $V_j\subseteq\E^2$ be the binary span after the first $j$ unrestricted coordinates. Both projections of $V_5$ have dimension $32$ except with probability $2B$, with $B$ defined above. On that event $V_5$ itself has dimension $32$ and both projections are injective. A nonzero intersection with a random diagonal multiple then requires a specified ratio in \emph{both} coordinates for some pair of nonzero vectors. Thus the probability that $\dim V_6<64$, on this event, is at most
\[
 J_6=\frac{(2^{32}-1)^2}{(Q-1)^2}.
\]

For a fixed $a$-dimensional projected subspace and $a\le b<2a$, the intersection-counting argument of~\eqref{eq:two-dimension-tail} gives
\[
 p(a,b)=\frac{(2^a-1)^2}{(Q-1)(2^{2a-b}-1)}
\]
as an upper bound on the probability that its next doubled dimension is at most $b$. On the event $\dim V_6=64$, an axis-supported subspace of $V_6$ has dimension $64-D$, where $D$ is the other projected dimension. Counting possible intersection pairs separately on the two axes and off the axes, then averaging over the sixth projected dimensions, bounds the additional failure to reach $\dim V_7=128$ by
\[
 J_7=\frac{(2^{64}-1)^2}{(Q-1)^2}
 +2\sum_{b=32}^{63}p(32,b)\frac{(2^{64-b}-1)^2}{Q-1}.
\]
These are bounds on intersected good-event probabilities; no conditioning denominator is omitted. In particular, the projected dimension tails are averaged before restricting to joint-dimension success.

The remaining four coordinates supply four independent random diagonal multiples of $V_7$, alongside $V_7$ itself. When $\dim V_7=128$, the two-coordinate annihilator bound is at most $2^{385-5\cdot128}=2^{-255}$, using $(Q/(Q-1))^8<2$. For a singleton coordinate support with projected dimension $D$, the corresponding bound is
\[
 h(D)=\min\{1,2^{193-5D}\}.
\]
We must average this expression, not discard a rare projected rank loss. Put
\[
 A(a)=h(2a)+\sum_{b=a}^{2a-1}p(a,b)h(b),
 \qquad A=A(64)+\sum_{a=32}^{63}p(32,a)A(a).
\]
The first expression bounds the seventh-step average given sixth projected dimension $a$; the second bounds its average after starting with dimension $32$ at step five. The two singleton-support errors therefore sum to at most $2A$. The total normalized error is at most $2B+J_6+J_7+2A+2^{-255}$.

Returning to equality weights costs $22/Q+44/(Q-1)$ by the preceding normalization argument on the $22$ coordinates of the two points. Geometric multipliers cost at most another $44/Q$. Exact rational arithmetic verifies the common bound
\[
 \zkCompactError:={}
 2B+J_6+J_7+2A+2^{-255}
 +\frac{66}{Q}+\frac{44}{Q-1}<2^{-158}.
\]
The ordinary-equality bound omits $44/Q$ and is smaller. All estimates are uniform in the fixed subspaces encountered along the construction.
\end{proof}

\paragraph{Consequence for Flock.} The first $96$ alternatives suffice for the feature ranks in \noteref{zk-flock-endpoints}, using $96\cdot640=61440$ independent bits in a public batch of $2^{18}$, rather than $2^{21}$. Reserving those $96$ complete banks leaves $65536$ arbitrary real-compression slots. The packed Flock column has log size $18+8=26$. The former batch required a packed column of log size $29$, already exceeding the current stack cap $28$ before any other column was placed. The audit checks an accepted reference layout with memory log $25$, non-BLAKE table logs $(19,19,19,19,20)$ and BLAKE log $18$. This verifies size compatibility for a candidate profile, not completion of all padding libraries in that profile.

\subsubsection{A random evaluation after fixed binary disclosures}

An initial PCS query lives in $\K$, not at a uniform extension-field point. The following lemma handles its fixed linear disclosure directly. It does not assume a second honest random point or ask for fresh source bits.

\begin{theorem}[Fixed-disclosure span]\label{thm:fixed-observation-span}
Let $S\subset\cube{12}$ be the $448$-position support with six arbitrary bits and at most one of the last six bits set. Fix any surjective binary linear map $\zkFixedMap:\Ftwo^S\to\Ftwo^b$, independently of an honest uniform $r\in\E^{12}$. Except with probability at most $\zkFixedError{b}$, the map
\[
 \zkSwapCoins\longmapsto
 \left(\sum_{s\in S}\zkSwapCoins_s\eq(r,s),\,
       \zkFixedMap(\zkSwapCoins)\right)
\]
is onto $\E\times\Ftwo^b$. For $Q=2^{192}$, the first-five error $B$ and sixth-step tails $p_d$ defined above, a uniform bound is
\[
 h_b(d)=\min\left\{1,\left(\frac Q{Q-1}\right)^6 2^{192+b-7d}\right\},
 \qquad
 \zkFixedError{b}=B+h_b(64)+\sum_{d=32}^{63}p_dh_b(d)+\frac{12}{Q}.
\]
In particular $\zkFixedError{64}<2^{-154}$. Uniform independent source bits therefore jointly hide an extension-field evaluation and any such fixed $64$-bit disclosure, including arbitrary fixed translations of both outputs.
\end{theorem}
\begin{proof}
Exclude coordinates $r_i=1$, costing at most $12/Q$. Normalize only the extension-field output by its nonzero common factor and set $t_i=r_i/(1+r_i)$, independently uniform in $\E\setminus\{1\}$. The first $64$ columns have extension coordinates $a_s=\prod_{i<6}t_i^{s_i}$; the six further blocks have coordinates $t_j a_s$ for $6\le j<12$. Their fixed binary coordinates are arbitrary columns of $\zkFixedMap$ and need no product structure.

Condition on the first six multipliers and let $d$ be the binary dimension of the span of the $a_s$. Under the trace pairing on $\E$ and ordinary binary pairing on $\Ftwo^b$, at most $2^{192+b-d}$ joint characters annihilate the first block. For a character with nonzero extension component, annihilating any further block imposes $d$ independent affine binary equations on its single multiplier $t_j$. There are at most $Q2^{-d}$ solutions in $\E$, hence probability at most $(Q/(Q-1))2^{-d}$ on the actual domain. The six multipliers are independent, giving $h_b(d)$ by a union bound. A character with zero extension component cannot annihilate all blocks because $\zkFixedMap$ is onto.

The extension projection reaches dimension $32$ in five steps except with probability $B$. After that, the sixth-step dimension lies in $[32,64]$ and has tails $p_d$. These bounds also apply on $\E\setminus\{1\}$: replacing multiplier $1$ by $0$ leaves $V+tV$ equal to $V$, so its subspace distribution is the same as on $\E^*$. Averaging $h_b(d)$ over the sixth dimension, rather than discarding every failed doubling, proves the formula. All estimates are uniform over the fixed map and intermediate subspaces.
\end{proof}

\paragraph{Conditional interpretation and limit.} On good $r$, every fiber $\zkFixedMap(\zkSwapCoins)=y$ leaves a uniform extension evaluation. Since the disclosed linear map is fixed independently of $r$ and onto, $y$ itself is uniform and independent of $r$ before conditioning. Thus the same average error bound holds for every fixed disclosed value. This does not allow an arbitrary map chosen after seeing $r$, or further disclosures of the same bits. The same formula gives only $\zkFixedError{128}<2^{-127}$ and $\zkFixedError{192}<2^{-63}$, without the required $128$-bit margin. Increasing the retained interface is not free in this lemma.

\paragraph{Evidence and application.} \auditref{zk_two_point_audit.py} computes the exact rational bounds and checks native joint ranks $256,320,384$ for fixed full-rank disclosures of $64,128,192$ bits. These samples check the implementation; the uniform proof is the character and dimension-tail argument above. \noteref{zk-flock-pair-opening} applies the $64$-bit case to the actual PCS weights and the same counter bits that hide a BLAKE2s terminal evaluation.

\subsubsection{A denser support for a larger retained interface}

Keep the same twelve coordinates, but take
\[
 \zkDenseSupport=\{s\in\cube{12}:s_0,\ldots,s_7\text{ arbitrary},\quad s_8+\cdots+s_{11}\le1\},
 \qquad |\zkDenseSupport|=1280.
\]
The inequality is over integers. This contains the preceding $448$-position support and fits the already allocated $4096$-pair BLAKE2s counter bank.

\begin{theorem}[Dense fixed-disclosure span]\label{thm:dense-fixed-observation-span}
Replace $S$ in Theorem~\ref{thm:fixed-observation-span} by $\zkDenseSupport$. For any fixed onto binary map to $\Ftwo^b$ with $b\le320$, the joint map fails to be onto with probability at most $\zkDenseError{b}<2^{-160}$. The same conclusion with error $\zkDenseError{b}+24/Q<2^{-160}$ holds when the extension weights are multiplied by any fixed nonzero geometric factors $\prod_i\lambda_i^{s_i}$.
\end{theorem}
\begin{proof}
After normalization, the first eight coordinates generate an extension-field projected subspace, and the last four supply independent scalar multiples of it. For any fixed projected dimension $d$, the preceding character argument gives
\[
 H_b(d)=\min\left\{1,\left(\frac Q{Q-1}\right)^4 2^{192+b-5d}\right\}.
\]
It remains to average three further doublings after reaching dimension $32$ at step five. For $a\le c<\min\{2a,192\}$ put
\[
 p(a,c)=\min\left\{1,\frac{(2^a-1)^2}{(Q-1)(2^{2a-c}-1)}\right\},
 \qquad m_a=\min\{2a,192\},
\]
and define
\[
 F_0(a)=H_b(a),\qquad
 F_{j+1}(a)=\min\left\{1,F_j(m_a)+\sum_{c=a}^{m_a-1}p(a,c)F_j(c)\right\}.
\]
The intersection-counting bound controls each exact dimension probability by the displayed tail. Nonnegativity proves inductively that $F_j(a)$ bounds the expected final error after $j$ more doublings, uniformly over the current subspace. Hence
\[
 \zkDenseError{b}=B+F_3(32)+12/Q.
\]
The same $0$-versus-$1$ substitution as above justifies the actual normalized multiplier distribution. Exact rational evaluation proves the stated bounds, uniformly for $b\le320$ by monotonicity in $b$.

For geometric weights, transform only the random extension point using $u_i=\lambda_i r_i/(1+r_i+\lambda_i r_i)$. Excluding its poles and coupling the resulting points to uniform costs at most $24/Q$, as in Lemma~\ref{lem:metadata-span}. The resulting common output factor is nonzero; the fixed disclosed map is unchanged.
\end{proof}

\paragraph{Budget and evidence.} The same formula gives $\zkDenseError{352}<2^{-158}$, but its value at $384$ exceeds $2^{-128}$. This is a limit of this estimate, not an impossibility theorem for the source. \auditref{zk_two_point_audit.py} reproduces the dimension recursion exactly and checks native joint ranks $384,448,512$ for $192,256,320$ retained bits. The use below retains one prior extension evaluation and one base-field query coordinate, requiring $256$ retained bits.
